\section{De «saber operar la computadora» a «entender la tarea digital»}

Las trayectorias de entrada de los tres invitados muestran precisamente que el GUI Agent no es una tendencia nueva que haya surgido de repente. Gu Yu (谷雨) empieza por el semantic parsing: la investigación temprana se centraba en cómo convertir una instruction en lenguaje natural en una consulta a una base de datos, una operación sobre una base de conocimiento u otra formal representation. En esa etapa el entorno era más controlable y la tarea se parecía más a «traducir una frase en una estructura ejecutable». Pero la naturaleza del problema ya se acercaba mucho a la del computer-using agent actual: el usuario da un objetivo en lenguaje natural y el sistema tiene que convertirlo en una acción ejecutable dentro del entorno.

La trayectoria de Xie Tianbao (谢天宝), en cambio, se despliega a lo largo de agent benchmark y entornos GUI reales. Desde AgentBench, intercode y OSWorld hasta el GUI grounding multimodal, lo que le interesa es cómo lograr que el modelo haga cosas de verdad en la pantalla, en páginas web y en el sistema operativo, y no solo responda preguntas en el espacio del texto. Shang Yanyi (尚晏仪) aporta una tercera línea: desde 2021 desarrollando productos capaces de operar la computadora en lugar del usuario, después haciendo Web Agent y GUI Agent de propósito general, hasta descubrir que una entrada de propósito general es demasiado difícil de convertir en producto comercial y virar hacia escenarios concretos como los viajes con IA. Al poner juntas las tres líneas, se ve que lo central del GUI Agent no es el «clic del ratón» en sí, sino cómo se conectan entre sí el lenguaje, la visión, la acción, las preferencias del usuario y los procesos de negocio reales.

\begin{knowledgebox}{Antecedentes históricos del GUI Agent}
\begin{tabularx}{\textwidth}{>{\raggedright\arraybackslash}p{0.26\textwidth} Y}
\toprule
Pista de la historia previa & Su significado para el GUI Agent actual \\
\midrule
Semantic Parsing & Entrenar sistemas para convertir tareas en lenguaje natural en representaciones ejecutables, lo que ofreció una forma temprana de «del lenguaje a la acción». \\
Web Automation / RPA & Demostró que muchas tareas digitales pueden convertirse en flujos de proceso, pero también expuso la fragilidad de los sistemas basados en reglas frente a entornos dinámicos. \\
Agent Benchmark & Usó tareas como OSWorld y Mind2Web para hacer explícito el problema de las capacidades, de modo que los modelos pudieran compararse entre sí. \\
Exploración de conversión en producto & Expuso la distancia entre un demo de investigación y la experiencia del usuario: poder ejecutar no equivale a ser útil, y poder mostrar el proceso no equivale a generar confianza. \\
\bottomrule
\end{tabularx}
\end{knowledgebox}

El cambio de 2025 es que estas líneas finalmente volvieron a conectarse gracias a modelos multimodales más potentes. Xie Tianbao menciona que los modelos tempranos ni siquiera podían realizar bien acciones aparentemente simples como «hacer clic en el botón de cerrar de la esquina superior derecha», «arrastrar el volumen a 80\%» o «modificar el nombre en el campo de entrada». Después, mediante datos sintéticos a gran escala, la expansión de entornos y la iteración de modelos, el GUI grounding pasó de ser «puro concepto» a resultar «básicamente utilizable». Esto no es un ajuste menor, porque la dificultad del entorno GUI no está solo en ver el botón, sino en alinear al mismo tiempo la localización visual, la instrucción en lenguaje, las coordenadas de la pantalla, el tipo de acción y el objetivo de la tarea.

Pero es precisamente aquí donde el problema del producto empieza a agudizarse. La actitud de Shang Yanyi al hablar de Operator es representativa: le hace sentir a quienes trabajan en el área que la dirección ha sido validada, pero ella no considera que sea un buen producto. Ver a una pequeña computadora hacer clic lentamente en una página web confunde con facilidad al usuario; una acción que él mismo podría completar en unos segundos, mostrada por la IA a una velocidad muy lenta, no necesariamente aumenta la confianza y puede, por el contrario, exponer torpeza. El investigador quiere demostrar que «de verdad operé por ti», mientras que el usuario tal vez solo quiere saber si «el resultado es confiable y si necesito confirmarlo».

\begin{warningbox}{Equiparar el GUI Agent con «ver la pantalla y hacer clic» hace perder de vista el problema real}
Las acciones superficiales del GUI Agent son hacer clic, escribir, desplazar y arrastrar; pero su problema profundo es la comprensión de la tarea digital. Un Agent maduro debería saber cuándo hacer clic en la GUI, cuándo llamar a una API, cuándo escribir un script, cuándo leer el archivo directamente y cuándo pedirle confirmación al usuario. La GUI es una parte del espacio de acción, no la inteligencia en sí misma.
\end{warningbox}

Lo que este capítulo establece en realidad es la premisa subyacente de toda la conversación: el futuro del GUI Agent no es solo un grounding más fuerte, ni tampoco una «pequeña computadora remota» más realista. Lo que tiene que resolver es la comprensión del objetivo, la selección de la acción, el aprendizaje a partir de la retroalimentación y los límites de la confianza dentro del entorno digital. Solo si esta premisa queda establecida, las discusiones posteriores sobre benchmark, modelo de negocio, UDA, memory y foso competitivo podrán conectarse en una sola línea.

\subsection{Resumen de la sección}

El GUI Agent no es un concepto nuevo que haya surgido de la nada, sino el resultado de la convergencia de semantic parsing, web automation, agent benchmark y la exploración de conversión en producto. La capacidad multimodal de 2025 hizo que «operar la computadora» entrara en un rango utilizable, pero lo verdaderamente difícil es pasar de «ejecutar la acción» a «entender la tarea digital». La controversia en torno a productos como Operator muestra que todavía hay una distancia considerable entre la demostración de investigación y la experiencia del usuario.
